\documentclass[11pt,a4paper]{article}
\usepackage[T1]{fontenc}
\usepackage{lmodern,amsmath,amssymb,amsthm,mathtools}
\usepackage[margin=25mm,headheight=14pt]{geometry}
\usepackage{microtype,enumitem,booktabs,tabularx,xcolor,fancyhdr}
\usepackage[colorlinks,linkcolor=blue!45!black,urlcolor=blue!45!black]{hyperref}
\hypersetup{pdftitle={A shorter route to RH and Dirichlet GRH: revised proof plan},
 pdfauthor={GRH distillation project}}
\setlength{\parskip}{4pt}
\setlength{\parindent}{0pt}
\setlength{\emergencystretch}{2em}
\setlist{topsep=3pt,itemsep=3pt,parsep=0pt}
\pagestyle{fancy}\fancyhf{}
\fancyhead[L]{\small RH and Dirichlet GRH}
\fancyhead[R]{\small Revised proof plan}
\fancyfoot[C]{\thepage}
\newtheorem{proposition}{Proposition}
\newtheorem{lemma}[proposition]{Lemma}
\DeclareMathOperator{\tr}{tr}
\DeclareMathOperator{\ord}{ord}
\DeclareMathOperator{\disc}{disc}
\DeclareMathOperator{\Rea}{Re}
\newcommand{\Q}{\mathbb Q}\newcommand{\C}{\mathbb C}
\newcommand{\R}{\mathbb R}\newcommand{\Z}{\mathbb Z}
\newcommand{\eps}{\varepsilon}
\newcommand{\repo}{https://github.com/JoseBrox/GRH_distillation/blob/main/}
\newcommand{\notelink}[2]{\href{\repo#1}{#2}}
\newcommand{\source}{https://github.com/lsernau/grh-lean/tree/72fead102f663d5c7a407964a13dc0cf9d977207}
\title{A shorter route to RH and Dirichlet GRH\\[4pt]
 \large A revised proof plan and its mathematical improvements}
\author{GRH distillation project}
\date{26 November 2026}
\begin{document}
\maketitle
\vspace{-1.3em}

\begin{abstract}
The preferred route now compares short integral forms with a lower bound
for the same determinant. This organizes the Sernau--Anthropic rewriting
of Alp\"oge's argument and incorporates stronger reductions proved in
our project. The analytic endpoint and several abstract tools are written
out; the common arithmetic construction remains to be reconstructed.
This document is a research plan, not a completed human proof of RH,
Dirichlet GRH, or general Hecke GRH.
\end{abstract}

\section{The idea to keep in view}

A nonzero algebraic number cannot be small at every place. Its ordinary
absolute values and its prime-adic absolute values satisfy the product
formula. The proof seeks two incompatible estimates for one carefully
chosen algebraic determinant.

Arithmetic divisibility makes normalized theta forms short. A second
moment shows that a small collection of their translates supplies almost
all the dimensions of the target. Separately chosen auxiliary twists
then force the corrected determinant of these forms to be large. The
large common term cancels. What remains is a bound for a prime sum whose
oscillations detect zeros of Dirichlet $L$-functions.

The square-root scale comes from a chain containing about $\sqrt{x}$
lifts. This is an explanation of the mechanism, not a substitute for
constructing those lifts or proving their norm, denominator and
nonvanishing estimates.

\subsection*{What the new source establishes at its public interface}

The pinned \href{\source}{source revision} exports ordinary RH and
\emph{Dirichlet} GRH for every positive modulus and every character,
including imprimitive characters. Its definitions are Mathlib's actual
analytically continued functions. The principal pole at $1$ is excluded;
the zero-on-the-line assertion concerns $0<\Rea s<1$. It does not export
Hecke GRH over arbitrary number fields.

This is a major reorganization of the same originating argument, not an
independent elementary method. The readable edition is the better base
for a conventional exposition. Its analytic error allowances must not
be mixed with the compact edition's different choices.

\newpage
\section{One determinant is the organizing invariant}

Let $D>0$ be the target dimension and let $y_1,\ldots,y_n$ be
independent selected images, with $0<n\le D$. Write $C$ for an upper
bound on their logarithmic \emph{squared} covolume, averaged over the
complex embeddings. Let $I$ be the sum of the local denominator costs
on both conjugate sides, and $Q$ the cost of completing these vectors
to a basis.

Selected twists supply a correction bound $B$. The determinant comparison
is
\begin{equation}\label{nonneg}
 0\le C+I+Q+B.
\end{equation}
Shortness and the corrected determinant ratio should give
\[
 C\le n(H-mR+E_A),\qquad B\le-DH+E_B.
\]
Here $H$ is the common main term, $m>0$ the number of lifts, and $R$
the exact signed prime remainder. Therefore
\begin{equation}\boxed{\quad
 nmR\le -(D-n)H+nE_A+E_B+I+Q.
 \quad}\label{master}
\end{equation}
This is the final arithmetic acceptance test. The sign of $H$ matters:
if $H\ge0$, the missing-rank term is favorable. Otherwise it needs its
own bound. It must not disappear through an informal replacement of
$n$ by $D$.

\subsection*{Why the scalar argument is short}

For a number field $K$ of degree $d$, use finite-place weights
$\omega_v=[K_v:\Q_p]/d$ and $|p|_v=p^{-1}$. Suppose one common
$g\in K^\times$ has, at every finite place, a factorization
$g=u_vr_v$. The local factors may vary with $v$. Suppose one common
$h\in K^\times$ satisfies
\[
 \log|u_v|_v\le a_v,\quad \log|hr_v|_v\le e_v,
 \quad \sum_v\omega_va_v=I,
\]
with finite sums and nonnegative $a_v,e_v$. If
\[
 \frac1d\sum_\sigma\log|\sigma g|\le C+Q,
 \qquad
 \frac1d\sum_\sigma\log|\sigma h|+
                  \sum_v\omega_ve_v\le B,
\]
then the product formulas for $g$ and $h$ give~\eqref{nonneg}.
The key point is that $g$ is fixed: one cannot apply the product formula
to a different global determinant at each prime.

For a Gram determinant, clearing the two conjugate sides costs twice
the symmetric lattice index. The completion cost can be its exact Schur
complement determinant rather than a worst norm times the missing rank.
Combining all correction factors before taking heights also preserves
cancellations. These refinements, their proofs and sharp obstruction
examples are in the \notelink{notes/lsernau-mathematical-extraction-2026-11-26.md}
{determinant note}.

\newpage
\section{Enough independent images from two moments}

The spanning problem need not require pointwise control of every map.
Let $A_t:S\to T$ be a finite family of maps between complex Hilbert
spaces, with probabilities $p_t$, and dimensions $d>0,D>0$. Define
\[
 K=\sum_t p_tA_tA_t^*,\qquad
 \delta=\frac1d\sum_t p_t\tr(I-A_t^*A_t)_+.
\]
The positive part replaces negative eigenvalues by zero, recording only
the loss of source norm. Assume $\eps\ge0$ and
\begin{equation}\label{moments}
 \tr K^2\le(1+\eps)d^2/D,\qquad e=2\delta+\eps<1.
\end{equation}

\begin{proposition}
Put $\lambda=d/(2D)$. Then $0<\lambda<1$. For every integer $k\ge0$,
the images of at most $k$
maps span a subspace whose missing share is at most
\[
 e+(1-e)(1-\lambda)^k.
\]
For $e<\eta<1$, a missing share at most $\eta$ therefore costs at most
\[
 k=\left\lceil\frac{2D}{d}
             \log\frac{1-e}{\eta-e}\right\rceil
 \quad\hbox{maps},\qquad
 kd\le d+2D\log\frac{1-e}{\eta-e}.
\]
\end{proposition}

\begin{proof}
Let $F\subset T$ have dimension $\alpha D$ and put $U_t=A_t^{-1}F$.
For any $U\subset S$,
$\dim U-\|A_tP_U\|_{\rm HS}^2\le\tr(I-A_t^*A_t)_+$.
Average this at $U_t$ and apply Hilbert--Schmidt Cauchy--Schwarz:
\[
 \sum_t p_t\dim U_t
 \le\delta d+\tr(P_FK)
 \le d\bigl(\delta+\sqrt{\alpha(1+\eps)}\bigr)
 \le\tfrac d2(1+\alpha+e).
\]
Some map adds at least $d(1-\alpha-e)/2$ dimensions. The missing share
$r$ thus obeys $r_{j+1}\le(1-\lambda)r_j+\lambda e$. Iterate from
$r_0=1$, stopping if $r\le e$. Finally,
$\delta d\ge d-D$ implies $d<2D$, and
$(1-\lambda)^k\le e^{-\lambda k}$ gives the displayed cost.
\end{proof}

For $A_g=\rho(g)B$ under a finite unitary group action, the moment is
\[
 \tr K^2=\frac1{|G|}\sum_{g\in G}
                         \|B^*\rho(g)B\|_{\rm HS}^2.
\]
Thus one scalar trace calculation suffices. A source restriction of
codimension $c$ loses at most $kc$ further dimensions. An orthogonal
target projection cannot increase the covariance moment, but its loss
of source norm must be counted again in $\delta$.

Choose independent vectors from the original short integral images.
Complex orthogonalization alone does not preserve their lattice. The
\notelink{notes/lsernau-moment-and-prime-reductions.md}{full proof}
also gives a sharper nonlinear recurrence and comparison with the
source's stronger uniform assumptions.

\newpage
\section{Twist specialization without circular choices}

The new route chooses an auxiliary anticyclotomic character at each
coefficient prime. One character must work simultaneously for a finite
family of representations and for all its required powers. The source
turns this into simultaneous specialization of polynomial minors of
finite presentations.

A potential circularity deserves explicit attention: the degrees of
those polynomials grow with the truncation level, while the torsion
character must already exist at that level. The following stronger
specialization statement removes degree dependence.

\begin{proposition}
Normalize the valuation on $\overline{\Q}_p$ by $v(p)=1$ and let
$\mathcal O$ be its valuation ring. Fix an integer $s_0>0$, $\rho>0$, $B\ge0$ and
$\eps>0$. There is an integer $H$, depending only on these parameters
and $p$, with the following property.

For any finite family of integral polynomials $P_i$ in at most $s_0$
variables, suppose there are possibly different witnesses $x_i$ with
coordinate valuations at least $\rho$ and
$\sum_i v(P_i(x_i))\le B$. There are roots of unity $\zeta_j$ of
$p$-power order dividing $p^H$, with
$0<v(\zeta_j-1)\le\rho$, such that
\[
 \sum_i\bigl(v(P_i(\zeta-1))-v(P_i(x_i))\bigr)_+\le\eps.
\]
In particular, $H$ is independent of all polynomial degrees and the
number of polynomials.
\end{proposition}

The proof uses the weighted Gauss valuation
$w_t(f)=\min_k(v(a_k)+kt)$ and its multiplicativity. Primitive
$p^h$-th roots lie on the shell of valuation
$r_h=1/\varphi(p^h)$. If $N_h$ roots of $f$ lie on that shell, then
$\sum_hN_hr_h\le B$. Since $\sum_h1/(h+2\delta)$ diverges,
one of finitely many shells has small enough root mass. Averaging
truncated distances on that shell gives a point whose valuation
excess is at most $\delta$. The truncation threshold $2\delta$
then shows that no distance was actually truncated.

Fix the common shell cutoff using the running budget $B+\eps$.
Replace the coordinates one at a time, using the product of the
current polynomials to share the error budget. Taking
$\delta=\eps/s_0$ proves the proposition. For a finite list of
nonzero signed powers $j$, the necessary radius loss is only
$p^{\max_jv_p(|j|)}$, rather than $\max_j|j|$. Negative powers
are cleared by factors that are units at all chosen points.

The \notelink{notes/lsernau-twist-specialization.md}{complete proof}
includes the shell bound, explicit choices, power substitutions and
edge cases. It chooses $H$ first and the cofinal presentation level
$N\ge H$ afterward. It does not construct the Selmer presentations,
their witnesses, or the comparison between polynomial valuations and
normalized module lengths. Those are the next arithmetic obligations.

\newpage
\section{The analytic endpoint is now explicit}

Fix $q\ge1$ and a primitive nonprincipal quadratic character $\kappa$
of conductor coprime to $q$. Let
\[
 a_\kappa(n)=\Lambda(n)(1-\kappa(n))\mathbf1_{n\equiv1\pmod q},
 \qquad \varphi=\varphi(q).
\]
For a principal character put $E_\psi(s)=(s-1)L(s,\psi)$, continued
at $1$; otherwise put $E_\psi=L(s,\psi)$. Define
\[
 P=\prod_{\psi\bmod q}E_\psi,\qquad
 Q_\kappa=\prod_{\psi\bmod q}L(s,\psi\kappa),\qquad
 G_\kappa=\frac1\varphi\left(\frac{Q_\kappa'}{Q_\kappa}
                       -\frac{P'}P-1\right).
\]
The constant $c_\kappa=G_\kappa(1)$ is real. The exact remainder is
\[
 R_\kappa(x)=\sum_{n\le x}a_\kappa(n)(1/n-1/x)
                  -\frac{\log x}{\varphi}-c_\kappa.
\]
Direct integration gives
\begin{equation}\label{transform}
 \int_1^\infty R_\kappa(x)x^{-s}\,dx
       =\frac{G_\kappa(s)/s-c_\kappa}{s-1}\qquad(\Rea s>1).
\end{equation}
The right side is regular at $1$. For $\Rea\rho>0$ and $\rho\ne1$,
its residue is
\[
 \frac{\ord_\rho P-\ord_\rho Q_\kappa}
                         {\varphi\rho(1-\rho)}.
\]

Suppose $1/2<\beta<1$ and $R_\kappa(x)\le Cx^{\beta-1}$ eventually.
If the real residues on $(\beta,1)$ are nonnegative, a positive-measure
Landau argument makes~\eqref{transform} holomorphic on $\Rea s>\beta$.
Indeed subtract it from $K/(s-\beta)$ to obtain the Laplace transform
of a nonnegative density. Such a transform cannot be analytic at its
finite abscissa of convergence; the residue sign excludes a pole there.
It follows that $P$ and $Q_\kappa$ have equal zero multiplicities
throughout that half-plane.

Given a proposed zero of $P$, it is enough to select a twist with no
zero at that point and no real zero in $(\beta,1)$. A stronger
finite-box selection theorem is now proved in ordinary mathematics:
for $8\mid q$, a unit $a\equiv7\pmod8$, and fixed $\beta,H$, there
are arbitrarily large primes $d\equiv a\pmod q$ such that every
$L(s,\psi\kappa_d)$ is nonzero for $\Rea s\ge\beta$, $|\Im s|\le H$.

Its proof uses a truncated M\"obius theta integral, three-lines,
a Jacobi-symbol second moment and quadratic reciprocity. Bad primes
up to $X$ number $O_{q,\beta,H}(X^{1-\delta/2})$, where
$\delta=(\beta-1/2)/9$. Their reciprocal sum converges, whereas
the reciprocal sum over the required progression diverges. All
normalizations and cancellation issues are checked in the
\notelink{notes/lsernau-zero-detection.md}{zero-detection proof}.

\newpage
\section{The six modules to reconstruct}

\textbf{1. Fix the exact prime remainder.}
For each putative zero, choose the accuracy and a finite rectangle,
then the quadratic twist, and finally let $x$ grow. The analytic
implication is now written out. Its only remaining nonclassical input
is the eventual one-sided arithmetic bound for that exact remainder.

\textbf{2. Construct one common theta diagram.}
Define the actual CM field, levels, Hermitian spaces, coefficient
lattices, measures, Poincar\'e map, target projection and two theta
chains. State the exact nonzero scalar in the commuting diagram.
Prove that the relevant target and determinants are nonzero. Every
reuse from our older construction needs a field, lattice and map
dictionary.

\textbf{3. Supply short integral images.}
Prove the theta-jet divisibilities and complex squared-norm bounds
for the same lattice. Keep the normalization, local clearing index,
ramified restrictions, auxiliary level and exponent rounding in one
ledger. A rescaling used for a moment calculation must retain its
effect on norms and denominators.

\textbf{4. Prove the two scalar moments.}
Establish~\eqref{moments} for the actual projected maps. Unfolding,
off-diagonal terms, multiplicities, stabilizers and masses remain
substantial. Choose the missing-rank tolerance jointly with the
completion and index costs; use the proved greedy lemma only afterward.

\textbf{5. Supply the selected-twist determinant bound.}
Construct finite Selmer presentations and the bounded witnesses needed
by specialization. Prove the integral congruence-to-projection-denominator
comparison, the one-lift determinant identity, and the chain/exchange
formula. One twist must serve the required finite family and powers at
a coefficient prime, including both signs in the critical-value ratios.
Combine all factors before assigning a summable all-prime budget.

\textbf{6. Assemble~\eqref{master}, RH first.}
Verify every bound at the same parameters and divide by the actual
positive rank. Begin with conductor one and the admissible baseline
$q_0=18480$, then carry an arbitrary fixed target conductor. Do not
put $q_0=1$ into a construction requiring $18480\mid q_0$.

\subsection*{Why the large cutoff is simpler than it first appears}

The readable source uses a fixed $q_0,d$ and a large cutoff $x_1$.
The classical prime number theorem in fixed progressions gives
$R_{q,d}(t)=O_{q,d,B}((\log t)^{-B})$ for every $B>0$, with the
constant above correctly normalized. If
$\log x_1\ge x^{\eps/12}$, these errors beat every prescribed power
of $x$. The source's exact comparison, $m\ge\sqrt{x}$ and
$\log b\ll x^\eps$ then give $R_{q_0,d}(x)\ll x^{-1/2+6\eps}$.
No moving-conductor prime number theorem is needed for this step.

\newpage
\section{What changes for Hecke GRH}

Dirichlet GRH gives the rational abelian portion of the Hecke problem.
For a finite-order Hecke character, its induced representation over
$\Q$ is a sum of one-dimensional characters exactly when the normal
closure of its class field is abelian over $\Q$. This precisely marks
the immediate Dirichlet reduction. The first cubic nonabelian benchmark
has an $S_3$ normal closure and lies beyond it.

The scalar product-formula argument works over number fields. The
source's selected-twist principle is also formulated over CM fields,
with explicit structural and arithmetic hypotheses. We can prove a
stronger structural input before attempting the geometric realization.

\subsection*{A useful class with zero structural defect}

Let $F$ be a CM field and $\lambda$ a finite odd-order anticyclotomic
character, nontrivial on inertia at some place away from an odd
coefficient prime $\ell$. Assume $\ell\nmid\ord(\lambda)\disc(F)$.
For finitely many rational finite-order characters $\eta_i$, assume
their orders are prime to $\ell$ and their least common multiple is
coprime to $\ord(\lambda)$. Let $\mathcal O$ be the integers of a finite
unramified extension of $\Q_\ell$ containing these finite character values.
Consider the lines
\[
 T_i=\mathcal O(\lambda^{\pm1}\eta_i\eps_\ell^{n_i}),
 \qquad n_i\in\Z,
\]
with strict (coboundary) local conditions above $\ell$ and their actual
ramified sets. The structural defects in the selected-twist theorem
are zero for every $T_i$.

Three elementary mechanisms account for the result. Coprime orders
let one choose common inertia killing every $\eta_i$ while retaining
a nontrivial odd-order scalar from $\lambda$; its square minus one is
a unit. At a genuinely ramified place away from $\ell$, prime-to-$\ell$
inertia kills the relevant dual invariants. Finally, any one-dimensional
residual rational character restricts to inertia at $\ell$ as a power
of the residual cyclotomic character. The dual ratio has exponent
$1-2n_i-2a$, which is odd and cannot vanish modulo $\ell-1$.
Since $F$ is unramified at $\ell$, restriction to $F$ preserves this
obstruction. No condition on the conductors of the $\eta_i$ is needed.

The \notelink{hecke/notes/lsernau-transfer-plan-2026-11-26.md}
{Hecke proof and exact interfaces} specialize this result to
$F=\Q(\sqrt{-3})$, the two nontrivial cubic sectors, all permitted
dyadic twists and every $\ell\ge5$. There is no high-dyadic threshold
in this structural theorem.

After the general twisting theorem is reconstructed, this gives
arbitrarily small \emph{normalized lengths after additional twists}.
It neither makes the original untwisted Selmer groups vanish nor
constructs the theta realization or determinant ratios. The exceptional
primes $2,3$, common fixed-CM geometry and critical-scale ratio estimate
remain. General infinite-order characters also retain their infinity
type, phase and archimedean-parameter obligations.

\newpage
\section{What to retain, and what to do next}

\begin{tabularx}{\textwidth}{@{}p{.31\textwidth}X@{}}
\toprule
Existing ingredient & Role in the revised route\\
\midrule
One-sided Mellin reasoning & Retain, with the exact new remainder and
all-character multiplicities checked.\\
Theta divisibilities, Poisson and scalar transport & Reuse after
identifying the actual spaces, levels, measures and maps.\\
Positive traces and integral restrictions & Supply or sharpen the two
scalar moments; keep projected source loss and lattice costs.\\
Smith windows, Fitting bounds and projectors & Improve the selected-twist
denominator comparison on its actual coefficient modules.\\
Original RH weighted Selmer result & Preserve in its established scope;
test compatibility before importing it into the new construction.\\
Geometric wall/divisor lower bound & Retain as an alternative supplier,
rather than a prerequisite of this determinant route.\\
General untwisted deformation propagation & Retain as an alternative;
prefer targeted selected-twist estimates for the new route.\\
\bottomrule
\end{tabularx}

\subsection*{The next decisive mathematical deliverables}

First reconstruct the finite-presentation and bounded-witness part of
the selected-twist theorem, then immediately apply it to the actual
integral denominator comparison. The specialization theorem is already
available, so this has a precise endpoint. In parallel construct the
common theta diagram and its scalar moments. Then prove the one-lift
and chain determinant identities and check the signed budget.

The old geometric and deformation branches should be expanded only
when they supply a specific needed lemma or resolve a concrete
obstruction. Our sampling and averaging criteria remain optional ways
to weaken the final arithmetic target; they are not extra tasks when
the stronger pointwise estimate is available.

\subsection*{Audit boundary and reading guide}

The \notelink{audit/lsernau-formal-audit-2026-11-26.md}{formal audit}
checks the public statements, pinned Mathlib definitions and source
inventory. \href{https://math.mit.edu/~drew/shortgrh/}{Andrew Sutherland's
independent verification} reports full builds, the three standard axioms,
all-module replay and successful sandboxed multi-kernel Comparator checks
for both editions at this exact revision. Michael Rubinstein also reports
an independent compact-edition pass. Our own local build is partial;
these external checks are attributed separately. No decisive flaw was
found in the inspected boundary and endgame. A conventional human audit
of all 803 readable modules remains beyond the work recorded here.

The \notelink{blueprint/merged-proof-plan.md}{canonical detailed plan}
and \notelink{audit/lsernau-audit-2026-11-26.md}{comparison audit}
give the exact source pins and dependencies. Every essential supplier
still needs its ordinary proof or a matching classical reference;
the source's \texttt{Lib} label alone is insufficient.

Existing map counts remain \textbf{RH 6/9}, \textbf{Dirichlet 18/31}
and \textbf{Hecke 6/9}. They count their original interfaces, not
percentages of this new route. A route change and its new abstract
lemmas do not close an old arithmetic node.

\end{document}
